\clearpage
\hypertarget{native-472-august}{}
\begin{center}
\rule{0.55\linewidth}{0.4pt}\\[2pt]
Astronomische Nachrichten. Band XXII. Nr. 513. S. 141..144. 22. August 1844.\\[2pt]
\rule{0.55\linewidth}{0.4pt}
\end{center}
\begin{flushright}Göttingen, 8. Aug. 1844.\end{flushright}

Seit dem 23. des vorigen Monats habe ich den Kometen zwar an mehreren Abenden kurz gesehen, doch erst gestern konnte ich wieder eine Beobachtung machen. Schon beim dritten Vergleich zogen aber erneut Wolken auf, die weitere Beobachtungen unmöglich machten. Ich erhielt die folgenden Werte, wobei ich zwei Sterne aus \textsc{Bessels} Zone Nr. 364 zum Vergleich heranzog.

\begin{center}
\setlength{\tabcolsep}{3pt}
\begin{tabular}{llll}
1844 Aug. 7. & $9^{\rm u}\ 16^{\rm m}\ 29^{\rm s}$ M. Z. & A. R. & $=213^\circ\ 52'\ 6''8$ \\
 & $9\ 24\ 48\mathord{\cdot}$\quad --- & Dekl. & $=+24\ 50\ 21.5$ \\
\end{tabular}
\end{center}

{\rightskip=0pt plus 1em Nach einer groben Überschlagsrechnung ergeben die parabolischen Elemente von Dr. \textsc{Goldschmidt} eine um $10''$ größere Rektaszension und eine um $16''$ kleinere Deklination.\par}

\begin{flushright}Göttingen, 10. Aug. 1844.\end{flushright}

Ich sende Ihnen erneut eine Positionsbestimmung des Kometen, die ich gestern Abend vorgenommen habe. Zum Vergleich dienten 2 Sterne aus \textsc{Bessels} Zone 364, der erste auch aus Zone 460.

\begin{center}
\setlength{\tabcolsep}{3pt}
\begin{tabular}{llll}
1844 Aug. 9. & $9^{\rm u}\ 36^{\rm m}\ 54^{\rm s}$ M. Z. & G. R. & $212^\circ\ 40'\ 23''8$ \\
 & $9\ 35\ 8$\quad --- & Abw. & $23\ 18\ 16.6$ \\
\end{tabular}
\end{center}

\textsc{Goldschmidts} Elemente ergeben eine um $5''$ größere Rektaszension und eine um $16''$ kleinere Deklination.
